\documentclass[conference]{IEEEtran}
\usepackage[T1]{fontenc}
\usepackage[utf8]{inputenc}
\usepackage{lmodern}
\usepackage{cite}
\usepackage{amsmath,amssymb,amsfonts}
\usepackage{algorithmic}
\usepackage{graphicx}
\usepackage{textcomp}
\usepackage{xcolor}
\usepackage{url}
\usepackage{booktabs}
\usepackage{multirow}

\usepackage[hidelinks]{hyperref}

\def\BibTeX{{\rm B\kern-.05em{\sc i\kern-.025em b}\kern-.08em
T\kern-.1667em\lower.7ex\hbox{E}\kern-.125emX}}

\begin{document}

\title{Detecção de Defeitos em Python via Modelos de Linguagem e Verificação Formal: Um Estudo Experimental}

\author{
\IEEEauthorblockN{Fernanda Alice Farias Duarte}
\IEEEauthorblockA{
\textit{Programa de Pós-Graduação em Informática, disciplina de Verificação Formal} \\
\textit{Universidade Federal do Amazonas} \\
Manaus, Brasil \\
fernanda.duarte@icomp.ufam.edu.br
}
}

\maketitle

\begin{abstract}
Detectar vulnerabilidades em código Python de forma confiável é difícil: modelos de linguagem são flexíveis, mas produzem falsos positivos; verificadores formais como o ESBMC-Python são precisos, mas custosos e limitados em escopo. Modelos de linguagem identificam padrões associados a defeitos com baixo custo, mas produzem alucinações e falsos positivos sem garantias formais, enquanto verificadores formais oferecem contraexemplos concretos, porém sem capacidade própria de triagem ou de distinção entre bugs executáveis e \textit{code smells}. Este trabalho propõe um pipeline híbrido composto por cinco etapas (Preprocess, LLM Analyzer, AST Validation, ESBMC e Report), no qual o modelo de linguagem gera hipóteses de defeitos, uma etapa de validação por AST filtra alucinações e apenas bugs verificáveis (divisão por zero, acesso fora dos limites e violação de asserção) são confirmados pelo ESBMC-Python, enquanto \textit{code smells} (método longo, excesso de parâmetros e condicional complexo) seguem direto para o relatório final. A proposta foi avaliada em um \textit{dataset} de 70 casos com \textit{ground truth} anotado manualmente, comparando três fluxos (somente ESBMC, somente modelo e híbrido) em cinco modelos de linguagem, sendo três de grande porte (LLMs) acessados via API e dois de pequeno porte (SLMs) executados localmente. Os resultados mostram que a confirmação formal reduz consistentemente os falsos positivos em todos os modelos, com ganhos mais expressivos para modelos locais de menor porte. A principal contribuição é demonstrar que combinar triagem heurística por modelo de linguagem com confirmação formal direta em Python aumenta a confiabilidade da detecção automática de defeitos e permite diferenciar, de forma sistemática, falhas executáveis de problemas de qualidade de código.
\end{abstract}

\begin{IEEEkeywords}
large language models, bounded model checking, formal verification, Python, software vulnerabilities, code smells
\end{IEEEkeywords}

\section{Introdução}

\subsection{Contexto}

Segundo o índice TIOBE de janeiro de 2026 \cite{tiobe2026}, Python ocupa a primeira posição entre as linguagens de programação mais utilizadas na atualidade. Sua popularidade se deve à simplicidade da linguagem e à ampla adoção em diferentes áreas, como desenvolvimento web, ciência de dados, inteligência artificial, entre outras. Como consequência, cresce também a necessidade de técnicas capazes de identificar defeitos, vulnerabilidades e problemas de qualidade em programas dessa linguagem.

Entretanto, características da própria linguagem Python, como a tipagem dinâmica, dificultam a análise automática do código e a verificação precisa de determinadas propriedades por técnicas tradicionais de análise estática \cite{farias2024}.


Nesse contexto, os Modelos de Linguagem para Código (\textit{Code Language Models}, Code LMs) vêm sendo utilizados em diversas tarefas de engenharia de software, incluindo geração e revisão de código, depuração, detecção de vulnerabilidades e análise da qualidade do código \cite{ullah2024,kaniewski2025,zhou2024}.


Em paralelo, surgiu o ESBMC-Python \cite{farias2024}, que tornou possível aplicar técnicas de verificação formal diretamente a programas Python, produzindo contraexemplos concretos quando uma violação é detectada.



\subsection{Problema}

Apesar dos avanços recentes, os modelos de linguagem para código ainda apresentam limitações quando são utilizados na análise automática de código. Seus resultados podem variar bastante de acordo com o modelo escolhido, a estratégia de \textit{prompting} e o contexto fornecido durante a análise, podendo produzir alucinações, inconsistências e falsos positivos \cite{ullah2024,tihanyi2025}. Como consequência, os resultados produzidos por esses modelos devem ser interpretados como hipóteses de defeitos ou vulnerabilidades, e não como confirmações de sua ocorrência, exigindo mecanismos adicionais de validação.




Avaliações experimentais confirmam essas limitações \cite{ding2025,ullah2024}, reforçando a necessidade de mecanismos adicionais capazes de validar ou refutar as suspeitas identificadas.


Nesse contexto, a verificação formal pode atuar como um mecanismo de confirmação para defeitos que possam ser expressos como propriedades verificáveis do programa. Entretanto, essa abordagem não se aplica diretamente a todos os problemas identificados durante uma revisão de código, especialmente aos \textit{code smells}, que representam problemas de qualidade e manutenibilidade, e não necessariamente violações executáveis de propriedades de segurança. Essa diferença torna necessária a distinção entre bugs executáveis, passíveis de confirmação formal, e \textit{code smells}, que devem ser analisados e reportados por outros critérios.


Embora sejam conceitos distintos, bugs executáveis e \textit{code smells} podem ser sinalizados por características semelhantes do código, como estruturas condicionais complexas, métodos extensos ou baixa clareza estrutural. Por essa razão, sua diferenciação pode ser difícil para abordagens baseadas exclusivamente em reconhecimento de padrões \cite{santos2023,souza2026}.



Embora trabalhos recentes indiquem que abordagens híbridas possam explorar a complementaridade entre modelos de linguagem para código e verificação formal \cite{tihanyi2025}, ainda existem poucos estudos que investigam essa integração especificamente para programas Python, sobretudo na validação de defeitos identificados por esses modelos e na distinção entre bugs executáveis e \textit{code smells}. Dessa forma, permanece em aberto a investigação de métodos capazes de aumentar a confiabilidade da análise automática de código em Python por meio da combinação dessas abordagens.

Diante desse cenário, surge a seguinte questão de estudo:

\textit{Como combinar modelos de linguagem para código e verificação formal para aumentar a confiabilidade da detecção automática de defeitos em programas Python, reduzindo falsos positivos e diferenciando bugs executáveis de \textit{code smells}?}




\subsection{Motivação e Justificativa}

O surgimento do ESBMC-Python tornou viável a aplicação de técnicas de verificação formal diretamente a programas Python, criando a oportunidade de investigar abordagens híbridas que combinem análise baseada em modelos de linguagem e verificação formal \cite{farias2024,tihanyi2025}.

Nessa abordagem, os modelos de linguagem atuam como geradores de hipóteses, sinalizando trechos suspeitos de conter bugs executáveis e \textit{code smells}. A verificação formal é então direcionada apenas às hipóteses de bugs passíveis de formalização, permitindo confirmar ou refutar essas suspeitas por meio de contraexemplos quando uma violação é encontrada \cite{shivaji2026}. Assim, a verificação formal é aplicada apenas aos trechos que o modelo identifica como possivelmente defeituosos, em vez de analisar todo o código.
 \cite{tihanyi2025}.


Além de investigar uma abordagem ainda pouco explorada para programas Python, este trabalho busca reduzir o esforço necessário para analisar os relatórios produzidos pelos modelos, aumentar a confiança nos bugs reportados e apoiar a diferenciação entre falhas executáveis e problemas de qualidade de código \cite{santos2023,souza2026}. Assim, a proposta explora a complementaridade entre a capacidade de triagem dos modelos de linguagem e o rigor da verificação formal.



\textbf{Objetivo Geral}

Desenvolver e avaliar um pipeline híbrido que combine modelos de linguagem e verificação formal para detectar bugs executáveis e \textit{code smells} em programas Python, aumentando a confiabilidade dos resultados e reduzindo falsos positivos.

\textbf{Objetivos Específicos}

\begin{itemize}

\item Construir um \textit{dataset} de funções Python com bugs executáveis, \textit{code smells} e casos limpos. Para cada função, será criado um arquivo JSON com os rótulos corretos, denominado \textit{ground truth}, indicando a presença ou ausência de defeitos e suas respectivas categorias.


\item Comparar cinco modelos de linguagem (LLMs e SLMs) na identificação de bugs verificáveis e \textit{code smells}, diferenciando-os de funções sem defeitos.

\item Fazer uma etapa de validação por AST para verificar se os problemas apontados pelo modelo realmente existem no código antes da confirmação formal.


\item Integrar o ESBMC-Python como confirmador formal das hipóteses de bug, restringindo sua execução apenas aos casos sinalizados pelo modelo como verificáveis.

\item Comparar o fluxo híbrido (Flow B) com a análise exclusiva por modelo de linguagem (Flow C), usando o ESBMC direto (Flow A) como instrumento de validação do \textit{dataset}.

\item Avaliar o impacto da confirmação formal na redução de falsos positivos, por meio das métricas de Precisão, Revocação, F1, TCF e TRN, detalhadas nas seções seguintes.

\end{itemize}


\section{Trabalhos Correlatos}

Nesta seção, são apresentados trabalhos relacionados ao uso de modelos de linguagem na análise de código, à verificação formal de programas Python e à combinação dessas duas abordagens. Também são discutidas as diferenças e relações entre bugs e \textit{code smells}.



\subsection{Modelos de Linguagem na Análise de Código: Capacidades e Limitações}

Embora modelos de linguagem tenham apresentado bons resultados em tarefas relacionadas a código, ainda enfrentam dificuldades na detecção de vulnerabilidades reais. O benchmark PrimeVul \cite{ding2025} mostrou que o desempenho desses modelos diminui quando avaliados em vulnerabilidades de códigos de projetos reais, atingindo F1-Score de apenas 3{,}09\%. Estudos também indicam que esses modelos podem ser influenciados por algumas características do código, como renomeações de variáveis, o que pode gerar achados inconsistentes e falsos positivos \cite{ullah2024}.

Outros estudos indicam que parte desse baixo desempenho pode estar relacionada à falta de contexto durante a análise \cite{li2026}. Quando informações adicionais sobre a execução e o fluxo de dados são fornecidas, modelos de raciocínio, como o DeepSeek-R1, apresentam resultados melhores, com F1 próximo de 0{,}6 em avaliações mais rigorosas. Ainda assim, esses modelos podem revisar uma análise inicialmente correta e chegar a uma conclusão incorreta após raciocínio excessivo.

Assim, embora modelos de linguagem sejam úteis para a triagem inicial e para identificar padrões no código, seus resultados devem ser tratados como hipóteses, e não como confirmações de defeitos.

\subsection{Verificação de Programas Python: Direta vs.\ Transpilação}

A verificação formal de programas Python tem sido tratada principalmente de duas formas. O ESBMC-Python \cite{farias2024} permite verificar programas Python diretamente, sem necessidade de tradução para outra linguagem de programação (detalhes na Seção~\ref{sec:esbmc-python}).

Por outro lado, sistemas como EVA \cite{shivaji2026} e PyVeritas \cite{orvalho2025} utilizam modelos de linguagem para traduzir programas Python para C e, em seguida, aplicam verificadores formais sobre o código traduzido. Essa estratégia permite reutilizar ferramentas de verificação já consolidadas para C, mas pode introduzir diferenças entre o comportamento do programa original e o código gerado.


A técnica de \textit{Bounded Model Checking} (BMC) \cite{biere1999}, utilizada pelo ESBMC, verifica propriedades dentro de um limite finito de execução. Para isso, o programa e a propriedade analisada são traduzidos em uma fórmula lógica resolvida por um \textit{solver} SMT.

\subsection{Integração entre Modelos de Linguagem e Verificação Formal}

O uso de modelos de linguagem junto à verificação formal tem sido investigado como forma de apoiar a análise automática de programas \cite{tihanyi2025}. Nessa abordagem, o modelo de linguagem pode gerar hipóteses, especificações, asserções ou invariantes, enquanto o verificador formal avalia se essas propriedades são satisfeitas pelo programa e produz contraexemplos quando encontra uma violação.

O sistema PropertyGPT, por exemplo, utiliza modelos de linguagem para gerar propriedades formais destinadas à análise de contratos inteligentes, isto é, programas executados automaticamente em plataformas de blockchain, como Ethereum. Dessa forma, o sistema busca identificar vulnerabilidades nesse tipo de aplicação \cite{tihanyi2025}.

O trabalho EVA \cite{shivaji2026} aplica essa ideia a programas Python por meio de um orquestrador que coordena ferramentas de análise e traduz Python para C antes da verificação. Seus resultados mostram que a verificação formal pode confirmar falhas que não são identificadas de forma confiável apenas pela análise do modelo de linguagem, especialmente em situações influenciadas pelas diferenças entre a semântica de Python e C.

\subsection{A Relação entre Bugs e \textit{Code Smells}}

A distinção entre falhas executáveis e problemas de qualidade de código é um desafio relevante na análise automática de programas. Santos et al.\ \cite{santos2023} identificaram associações entre bugs e \textit{code smells}, mostrando que essas categorias podem compartilhar características estruturais semelhantes. Essa sobreposição ajuda a explicar por que modelos de linguagem podem confundir falhas executáveis com problemas de manutenibilidade.

Estudos comparativos também mostram que modelos de linguagem conseguem identificar alguns \textit{smells} estruturais mais objetivos, como \textit{Long Method} e \textit{Large Class}, mas apresentam resultados mais variáveis em categorias que dependem de maior interpretação do contexto \cite{souza2026}.


\subsection{Posicionamento da Proposta}

Em síntese, o ESBMC-Python \cite{farias2024} fornece infraestrutura para a verificação formal direta de programas Python, mas, isoladamente, não realiza uma triagem inicial do código nem trata problemas de qualidade, como \textit{code smells}. Por outro lado, trabalhos que utilizam modelos de linguagem para detectar vulnerabilidades e \textit{smells} \cite{ding2025,li2026,souza2026} ampliam a capacidade de triagem, mas seus resultados permanecem heurísticos.

O trabalho EVA \cite{shivaji2026} é o mais próximo desta proposta, pois combina modelos de linguagem e verificação formal para programas Python. Contudo, depende da tradução do código Python para C antes da verificação e foi avaliado em um conjunto próprio de 23 programas.

A proposta deste trabalho utiliza modelos de linguagem para apontar possíveis problemas no código e distingue dois fluxos de tratamento. Bugs executáveis que podem ser expressos como propriedades verificáveis seguem para confirmação pelo ESBMC-Python. Já os \textit{code smells} não passam por confirmação formal, pois representam problemas de qualidade e manutenção do código. Nesses casos, o relatório mantém a identificação feita pelo modelo, junto com a explicação e uma possível sugestão de refatoração.


Além disso, a proposta inclui uma etapa de validação por AST antes da verificação formal, para filtrar achados que não correspondem à estrutura do código. O trabalho também compara três abordagens: análise apenas com modelos de linguagem, verificação formal isolada e o fluxo híbrido. Essa comparação utiliza Precisão, Revocação e F1-Score como métricas principais, complementadas por Acurácia, MCC, TCF e TRN, detalhadas na Seção~\ref{sec:metricas}. Por fim, o \textit{ground truth} é definido separadamente das saídas do pipeline, permitindo avaliar de forma mais clara se a confirmação formal contribui para reduzir falsos positivos.



\section{Referencial Teórico}

\subsection{Bounded Model Checking e Resolvedores SMT}

O \textit{Bounded Model Checking} (BMC) é uma técnica de verificação formal que explora o espaço de estados de um sistema até um limite finito $k$ \cite{biere1999}. O processo consiste em traduzir o programa e a negação da propriedade desejada em uma fórmula lógica. Se a fórmula for satisfatível (\textbf{SAT}), o resolvedor SMT (\textit{Satisfiability Modulo Theories}) produz um contraexemplo concreto que demonstra uma violação da propriedade. Caso seja insatisfatível (\textbf{UNSAT}), não foi encontrada violação dentro do limite analisado.

Os resolvedores SMT estendem a verificação booleana tradicional ao permitir o uso de teorias como aritmética, vetores de bits e arranjos. Essa característica torna o BMC útil para encontrar bugs que se manifestam em execuções curtas e violações em casos extremos. Entretanto, a ausência de violações até o limite $k$ não garante que o programa esteja livre de falhas em execuções mais longas.

\subsection{O Motor de Verificação ESBMC-Python}
\label{sec:esbmc-python}

O ESBMC \cite{cordeiro2012} é um verificador baseado em BMC e em resolvedores SMT, capaz de detectar automaticamente violações como divisão por zero, acesso fora dos limites e falhas de asserção, produzindo contraexemplos concretos. Sua extensão para Python, denominada ESBMC-Python \cite{farias2024}, permite verificar programas escritos diretamente nessa linguagem.

Para isso, o ESBMC-Python utiliza o \textit{parser} da linguagem para gerar a Árvore Sintática Abstrata (\textit{Abstract Syntax Tree}, AST) do programa. Em seguida, realiza inferência de tipos a partir das anotações e do uso das variáveis, traduzindo o código para a representação intermediária GOTO utilizada pelo ESBMC. A partir disso, o programa passa pela conversão para a forma SSA (\textit{Static Single Assignment}) e pela geração de uma fórmula lógica submetida a um \textit{solver} SMT. Essa fórmula pode ser representada como $C \land \neg P$, em que $C$ corresponde às restrições do programa e $P$ à propriedade de segurança analisada.

No presente trabalho, o ESBMC-Python é utilizado como etapa de confirmação formal. Apenas hipóteses de bugs classificadas como verificáveis (\texttt{verifiable=true}) pelo modelo de linguagem e validadas pela etapa de AST são encaminhadas para verificação formal.

\subsection{Modelos de Linguagem para Código}

Modelos de linguagem para código são redes neurais baseadas na arquitetura \textit{Transformer}, treinadas sobre grandes volumes de texto e código-fonte. Esses modelos podem interpretar e gerar código em diferentes linguagens de programação \cite{dakhel2023}. No contexto da análise de software, podem apoiar a triagem de bases de código e identificar padrões associados a vulnerabilidades ou \textit{code smells}.

Entretanto, seus resultados são heurísticos e não oferecem garantias formais sobre o comportamento do programa. Por isso, podem produzir alucinações, isto é, afirmações sobre o código que não correspondem à sua estrutura ou comportamento real \cite{ullah2024,ding2025}.

Neste trabalho, são avaliados modelos de grande porte (\textit{Large Language Models}, LLMs) acessados via API e modelos de pequeno porte (\textit{Small Language Models}, SLMs) executados localmente via Ollama. Essa escolha permite comparar diferentes perfis de custo, latência e qualidade de triagem.

\subsection{Bugs, Vulnerabilidades e \textit{Code Smells}}

Para fins deste trabalho, adota-se uma distinção entre bugs, vulnerabilidades e \textit{code smells}. Um \textit{bug} é um defeito que se manifesta como violação do comportamento esperado durante a execução, como uma divisão por zero, um acesso fora dos limites de uma estrutura de dados ou a violação de uma asserção explícita.

Uma vulnerabilidade corresponde a um bug, ou a uma classe de bugs, cuja exploração pode comprometer propriedades de segurança do sistema, como confidencialidade, integridade ou disponibilidade.

Um \textit{code smell} é um indicador de possível problema de qualidade ou manutenibilidade que não representa necessariamente uma falha executável imediata \cite{souza2026}. Exemplos incluem métodos excessivamente longos (\textit{long method}), listas de parâmetros extensas (\textit{many parameters}) e condicionais com complexidade lógica elevada (\textit{complex conditional}).




\section{Método Proposto}

\subsection{Visão Geral do Pipeline}
\label{sec:pipeline-overview}

O pipeline tem cinco etapas sequenciais: \textit{Preprocess}, \textit{LLM Analyzer}, \textit{AST Validation}, \textit{ESBMC} e \textit{Report}.

\textbf{Preprocess:} cada arquivo Python do dataset é carregado e processado individualmente. O módulo \texttt{ast} da linguagem converte o código em uma Árvore Sintática Abstrata (AST), que representa a estrutura do programa de forma organizada: quais funções existem, quais variáveis são usadas, quais operações são realizadas e onde cada elemento aparece. A AST extraída nesta etapa é usada internamente na etapa de AST Validation para checar as hipóteses do modelo, mas não é enviada ao modelo durante a análise. Pois incluir a AST no \textit{prompt} pode revelar ao modelo a estrutura exata de operações do código, introduzindo viés nos resultados.

\textbf{LLM Analyzer:} o código-fonte e metadados básicos da função (nome, parâmetros, anotações de tipo, número de linhas) são enviados ao modelo via \textit{prompt} estruturado com \textit{Chain-of-Thought} (CoT), técnica em que o modelo é instruído a raciocinar passo a passo antes de emitir um veredito, o que reduz erros em tarefas de análise \cite{li2026}. O \textit{prompt} do sistema define as seis categorias aceitas pelo pipeline (\texttt{division\_by\_zero}, \texttt{out\_of\_bounds} e \texttt{assertion\_violation} para bugs; \texttt{long\_method}, \texttt{many\_parameters} e \texttt{complex\_conditional} para \textit{code smells}, detalhadas na Seção~\ref{sec:dataset}) . O modelo é instruído a retornar exclusivamente um JSON válido com um schema fixo: para cada hipótese, o JSON inclui o tipo do achado (\texttt{suspected\_bug} ou \texttt{smell\_heuristic}), a categoria, a flag \texttt{verifiable}, uma explicação do raciocínio e a expressão exata identificada no código. A saída é processada por um \textit{parser} estruturado com extração via regex como alternativa, garantindo que respostas malformadas não interrompam o pipeline.

\textbf{AST Validation:} Verifica se a expressão reportada pelo modelo realmente existe no código da função analisada. A validação ocorre em três fases. Na primeira, a expressão é comparada diretamente contra as operações extraídas pela etapa de Preprocess (subscripts para \texttt{out\_of\_bounds}, divisões para \texttt{division\_by\_zero}): se houver correspondência exata, a hipótese é confirmada e recebe os metadados de linha. Na segunda fase, caso a expressão não tenha sido encontrada na fase anterior, ela é normalizada via AST (\texttt{ast.unparse}) e buscada nos nós executáveis do código-fonte; isso tolera diferenças superficiais de formatação, mas ainda exige que a expressão exista literalmente no código. Para \texttt{assertion\_violation}, a validação compara a expressão reportada com as instruções \texttt{assert} presentes na função, por correspondência de texto ou de nós AST. Se nenhuma das fases confirmar a expressão, a hipótese é marcada como alucinação (\texttt{llm\_false\_positive}) e descartada, evitando que o ESBMC seja invocado sobre suposições sem base no código real.

\textbf{Bifurcação:} após a AST Validation, o pipeline se divide em dois caminhos. \textit{Code smells} e bugs com \texttt{verifiable=false} seguem diretamente para o \textit{Report} com o veredito e a recomendação do modelo. Isso evita invocar o ESBMC sobre problemas que não têm natureza verificável formalmente, reduzindo o custo computacional. Apenas bugs com \texttt{verifiable=true} e validados pela AST seguem para a etapa de verificação formal.

\textbf{ESBMC:} o verificador é executado diretamente sobre o arquivo Python original, sem nenhuma modificação no código. A flag \texttt{{-}{-}function} indica ao ESBMC qual função verificar, com base na função identificada pelo modelo. A flag \texttt{{-}{-}assign-param-nondet} faz com que todos os parâmetros de entrada sejam tratados como valores desconhecidos: o ESBMC explora simbolicamente todas as combinações possíveis de entrada, sem precisar de casos de teste escritos manualmente. A flag \texttt{{-}{-}incremental-bmc} ativa o modo de verificação incremental, em que o limite de execução começa em um valor pequeno e aumenta gradualmente a cada iteração, permitindo encontrar violações com o menor limite possível e tornando a verificação mais eficiente do que fixar um valor único desde o início. As propriedades verificadas, divisão por zero, acesso fora dos limites de listas e violação de asserções, são checadas nativamente pelo verificador. Se encontrar uma situação que viola alguma dessas propriedades, o ESBMC retorna \texttt{VERIFICATION FAILED} junto com um contraexemplo concreto mostrando os valores que disparam o bug. Caso não encontre violação, retorna \texttt{VERIFICATION SUCCESSFUL}. A chamada gerada pelo pipeline para cada hipótese de bug segue o formato:

\begin{verbatim}
esbmc --function <nome_da_funcao> \
      --incremental-bmc \
      --assign-param-nondet \
      <arquivo.py>
\end{verbatim}

\textbf{Report:} consolida, por função, o veredito do modelo, o resultado do ESBMC (quando aplicável) e o rótulo do \textit{ground truth}, calculando as métricas por categoria e por modelo para os três fluxos. Os resultados são exportados em JSON e podem ser visualizados em um \textit{frontend} HTML estático gerado pelo pipeline.

\subsection{Fluxos de Avaliação}

A verificação formal via BMC pode demandar alto custo computacional e ser difícil de aplicar em bases de código maiores. À medida que a complexidade do programa cresce, o número de estados a serem analisados também aumenta, elevando o tempo necessário para a verificação \cite{tihanyi2025,cordeiro2012}. Além disso, em projetos reais, a aplicação da verificação formal pode exigir um trabalho prévio de preparação. Isso envolve selecionar e preparar funções ou trechos específicos do código, além de definir as propriedades ou asserções que devem ser verificadas. Esse trabalho torna mais difícil analisar diretamente todos os arquivos de um repositório grande. Shivaji et al.\ \cite{shivaji2026} reportam tempos de 35--50 segundos por função para a verificação formal, enquanto uma triagem por modelo de linguagem pode ocorrer em 2--10 segundos. Dessa forma, utilizar modelos de linguagem como etapa inicial permite direcionar o ESBMC apenas para funções ou trechos com indícios de bugs verificáveis, em vez de aplicar a verificação formal sobre todo o código. Essa escolha reduz o número de casos que precisam ser preparados e analisados, tornando o uso da verificação formal mais viável na prática \cite{kaniewski2025}.


O pipeline tem três modos de execução, chamados de fluxos:

\textbf{Flow A (ESBMC direto):} o ESBMC é executado diretamente sobre cada função, sem nenhuma triagem prévia por modelo de linguagem. Neste trabalho, esse fluxo foi utilizado para validar o \textit{dataset}, confirmando que os bugs anotados no \textit{ground truth} são de fato disparáveis pelo verificador. 


\textbf{Flow B (Híbrido):} o modelo faz a triagem inicial, a AST Validation filtra alucinações, e o ESBMC confirma ou refuta cada hipótese de bug verificável. É o fluxo principal da proposta.

\textbf{Flow C (Somente Modelo):} usa apenas o veredito do modelo de linguagem, após a AST Validation, sem confirmação formal. Serve para isolar o ganho específico da etapa de verificação.

A comparação principal é entre Flow~C e Flow~B: Flow~C mostra o que o modelo reporta sozinho; Flow~B mostra o efeito da combinação com o ESBMC, principalmente a redução de falsos positivos.

\subsection{Dataset e Categorias Avaliadas}
\label{sec:dataset}

O conjunto de dados (V1) é composto por 70 casos, cada um correspondendo a uma \textbf{função Python isolada}. Os casos são distribuídos entre 45 funções com \textit{bugs} executáveis (15 por categoria de bug), 10 controles limpos (funções corretas, sem defeitos, utilizadas para avaliar a taxa de falsos positivos) e 15 funções com \textit{code smells} (5 por categoria). É importante notar que o \textit{ground truth} registra instâncias de bug individualmente: três funções da categoria \texttt{out\_of\_bounds} contêm dois acessos inválidos distintos cada uma, totalizando 18 instâncias nessa categoria e 48 instâncias de bug no total. Por isso, nas tabelas de resultados, o número de verdadeiros positivos (VP) pode superar 15 nessa categoria — cada instância corretamente identificada conta como um VP separado.

São consideradas três categorias de \textit{bugs} verificáveis por ESBMC: \texttt{division\_by\_zero} (divisão por zero), \texttt{out\_of\_bounds} (acesso fora dos limites de uma estrutura indexável) e \texttt{assertion\_violation} (violação de uma asserção explícita no código). São consideradas também três categorias heurísticas de \textit{code smells}, definidas por critérios estruturais objetivos: \texttt{long\_method} (métodos com 20 ou mais linhas executáveis), \texttt{many\_parameters} (funções com 5 ou mais parâmetros, excluindo \texttt{self}/\texttt{cls}) e \texttt{complex\_conditional} (condicionais com 3 ou mais operadores lógicos). A avaliação por categoria, além da avaliação global, permite identificar quais tipos de defeito são mais ou menos acessíveis a cada modelo, informação relevante para guiar a seleção de modelos em cenários reais.

O \textit{ground truth} é o conjunto de rótulos corretos definidos manualmente para cada caso do dataset, indicando o tipo real de defeito presente (ou a ausência de defeito). 

Um passo crítico da construção do dataset é a \textbf{validação pelo ESBMC}: cada função anotada como contendo um bug verificável é submetida ao ESBMC-Python para confirmar que o defeito é disparável. Todos os 45 arquivos de bug foram confirmados com \texttt{VERIFICATION FAILED}, com o tipo de violação correspondendo à categoria anotada no \textit{ground truth}. 

\subsection{Modelos e Configurações Experimentais}

Cinco modelos foram selecionados para representar diferentes perfis de uso: \texttt{claude-sonnet-4-6} (Anthropic, via API), \texttt{gemini-3.1-flash-lite} (Google, via API), \texttt{gpt-5.5} (OpenAI, via API), \texttt{deepseek-r1:7b} e \texttt{qwen2.5-coder:7b} (modelos de pequeno porte e código aberto, executados localmente via Ollama). O mesmo \textit{prompt} foi utilizado para todos os modelos, sem ajustes específicos por modelo, caracterizando uma avaliação \textit{cross-model}. A temperatura foi configurada como 0 para Gemini, DeepSeek e Qwen. Para Claude e GPT-5.5, esse parâmetro não estava disponível nas respectivas APIs.

Os modelos locais (DeepSeek e Qwen) foram executados em uma máquina com processador Intel Core i7-12700H (12ª geração), 16\,GB de RAM e sem GPU dedicada, utilizando inferência apenas por CPU via Ollama.


\subsection{Métricas de Avaliação}
\label{sec:metricas}

Para cada função, o modelo pode gerar uma ou mais suspeitas de defeito. As métricas são calculadas em dois níveis, que respondem a perguntas diferentes:

\begin{itemize}
\item \textbf{Por suspeita:} cada suspeita reportada pelo modelo é avaliada individualmente. Esse nível permite medir quantas suspeitas estavam corretas e quantas eram falsos positivos. Por exemplo, se o modelo reporta duas suspeitas para uma função e apenas uma corresponde ao \textit{ground truth}, são contabilizados um VP e um FP.

\item \textbf{Por função:} cada função é avaliada como um todo: o modelo acertou ou errou ao classificar aquela função como bugada ou limpa, sem considerar categoria ou expressão específica. Esse nível é necessário para calcular métricas que dependem de verdadeiros negativos (TN): quando o modelo não reporta defeitos em uma função limpa, esse resultado é considerado um acerto.
\end{itemize}

No nível de suspeita, não há uma forma direta de contabilizar verdadeiros negativos, pois não é possível enumerar todas as suspeitas que o modelo deixou de reportar corretamente. Por isso, Precisão, Revocação e F1-Score são calculadas por suspeita. Já Acurácia e MCC, que dependem de TN em suas fórmulas, são calculadas por função.

Os resultados possíveis são:

\begin{itemize}
\item \textbf{Verdadeiro Positivo (VP)}: o modelo reportou uma suspeita que corresponde a um defeito presente no \textit{ground truth}.
\item \textbf{Falso Positivo (FP)}: o modelo reportou uma suspeita que não corresponde a nenhum defeito presente no \textit{ground truth}. Hipóteses descartadas pela AST Validation como alucinações também são contabilizadas como FP.
\item \textbf{Falso Negativo (FN)}: existe um defeito no \textit{ground truth}, mas o modelo não o reportou.
\item \textbf{Verdadeiro Negativo (TN)}: a função está limpa e o modelo não reportou defeitos. Contabilizada apenas no nível de função e apenas para bugs — o dataset não contém funções rotuladas como ``sem \textit{smell}'', portanto TN não se aplica a \textit{code smells}.
\item \textbf{Alucinação (Aluc.)}: hipótese reportada pelo modelo que referencia uma variável ou operação inexistente no código, detectada pela AST Validation. Aplicável apenas a bugs, pois a AST Validation não verifica hipóteses de \textit{code smells}. Contabilizada separadamente como diagnóstico da qualidade do modelo, mas também incluída no FP para fins de cálculo de precisão.
\end{itemize}


A escolha das métricas é fundamentada pela literatura e pelas necessidades específicas do pipeline híbrido. Precisão, Revocação e F1-Score são o padrão em benchmarks de detecção de defeitos com modelos de linguagem \cite{ding2025,li2026}. Acurácia e MCC são adotadas porque o \textit{dataset} é desbalanceado (mais funções com bugs do que limpas): Ding et al.\ \cite{ding2025} alertam que a acurácia isolada pode ser enganosa nesse cenário, enquanto o MCC considera todos os quadrantes da matriz de confusão e é mais robusto. A TCF e a TRN são métricas específicas para pipelines híbridos, onde a verificação formal atua como árbitro das hipóteses heurísticas \cite{tihanyi2025,shivaji2026}; a TRN em particular quantifica a redução de falsos positivos, lacuna identificada na literatura de refutação de bugs espúrios via SMT \cite{gadelha2019smt}. A separação em dois níveis (achado e função) e a exclusão de \textit{smells} das métricas formais derivam da distinção estrutural entre bugs verificáveis e problemas de qualidade \cite{santos2023}.

As três métricas a seguir se aplicam tanto a bugs quanto a \textit{code smells}, calculadas ao nível de achado.

A \textbf{Precisão} ($P$) responde: dos defeitos que o modelo acusou, quantos realmente existiam?
\begin{equation}
P = \frac{VP}{VP + FP}
\end{equation}
Alta precisão significa poucos alarmes falsos. Baixa precisão significa que o modelo acusa muita coisa que não é defeito.

A \textbf{Revocação} ($R$) responde: dos defeitos reais existentes no dataset, quantos o modelo encontrou?
\begin{equation}
R = \frac{VP}{VP + FN}
\end{equation}
Alta revocação significa que o modelo encontrou a maioria dos bugs. Baixa revocação significa que perdeu muitos.

O \textbf{F1-Score} combina precisão e revocação em um único valor, penalizando quando qualquer uma das duas é baixa \cite{ding2025}:
\begin{equation}
F1 = 2 \cdot \frac{P \cdot R}{P + R}
\end{equation}
É útil para comparar modelos com diferentes trocas entre precisão e revocação.

As quatro métricas a seguir são específicas para bugs. Para \textit{code smells}, o ESBMC não é aplicado, portanto TCF e TRN não são calculadas; Acc e MCC também não se aplicam a smells porque o dataset não tem funções limpas rotuladas como ``sem smell''.

A \textbf{Taxa de Confirmação Formal (TCF)} mede a qualidade das hipóteses que o modelo envia ao ESBMC. A LLM encaminha tanto bugs reais quanto falsos positivos para verificação; o ESBMC confirma os reais (\texttt{VERIFICATION FAILED}) e rejeita os falsos (\texttt{VERIFICATION SUCCESSFUL}). Sendo $C$ o número de hipóteses confirmadas e $N$ o número rejeitadas ou inconclusivas:
\begin{equation}
TCF = \frac{C}{C + N}
\end{equation}
Alta TCF indica que o modelo encaminha principalmente hipóteses reais. Baixa TCF indica que envia muitos falsos positivos ao verificador.

A \textbf{Taxa de Redução de Ruído (TRN)} responde: o ESBMC eliminou qual proporção dos falsos positivos que o modelo tinha sozinho?
\begin{equation}
TRN = \frac{FP_C - FP_B}{FP_C}
\end{equation}
onde $FP_C$ e $FP_B$ são os falsos positivos nos Flows C e B. TRN\,=\,1{,}000 indica que o ESBMC eliminou todos os FPs do modelo; TRN\,=\,0 indica que não eliminou nenhum. Quando $FP_C = 0$, a TRN é indefinida.

A \textbf{Acurácia por Função (Acc)} e o \textbf{Coeficiente de Correlação de Matthews (MCC)} são calculados ao nível de função, onde cada função é classificada como positiva (contém bug) ou negativa (limpa). O subscrito $f$ indica contagens por função.

A Acurácia responde: qual proporção de funções foi classificada corretamente?
\begin{equation}
\text{Acc} = \frac{VP_f + TN_f}{VP_f + FP_f + FN_f + TN_f}
\end{equation}
Alta acurácia significa que o modelo acertou a maioria das funções. Porém, como o dataset tem mais funções com bugs (45) do que limpas (10), a acurácia pode ser enganosa: um modelo que acusa bug em tudo teria acurácia de 82\% sem ser útil.

O MCC corrige esse problema avaliando o desempenho mesmo com classes desbalanceadas:
\begin{equation}
\text{MCC} = \frac{VP_f \cdot TN_f - FP_f \cdot FN_f}{\sqrt{(VP_f+FP_f)(VP_f+FN_f)(TN_f+FP_f)(TN_f+FN_f)}}
\end{equation}
Seu valor vai de $-1$ a $+1$: $+1$ é classificação perfeita, $0$ equivale a chutar aleatoriamente, e valores negativos indicam que o modelo é pior que o aleatório.

\section{Resultados}

O benchmark V1 foi executado sobre 70 funções Python (45 com bugs, 15 com \textit{code smells} e 10 limpas) com cinco modelos divididos em dois grupos:

\begin{itemize}
\item \textbf{Modelos de API} (\texttt{gpt-5.5}, \texttt{claude-sonnet-4-6}, \texttt{gemini-3.1-flash-lite}): modelos grandes acessados remotamente via API, com maior capacidade de raciocínio.
\item \textbf{Modelos locais} (\texttt{deepseek-r1:7b}, \texttt{qwen2.5-coder:7b}): modelos menores executados localmente via Ollama, sem custo de API mas com menor capacidade.
\end{itemize}

Para cada modelo são reportadas métricas no Flow~C (somente modelo de linguagem, sem verificação formal) e no Flow~B (híbrido: modelo + ESBMC). O Flow~A (ESBMC direto, sem modelo) foi executado separadamente para validação do \textit{dataset}, confirmando todos os 45 arquivos de bug com \texttt{VERIFICATION FAILED} (P\,=\,1{,}000, R\,=\,1{,}000, F1\,=\,1{,}000).

\subsection{Resultados Globais por Modelo}

A Tabela~\ref{tab:bugs-llm} apresenta as métricas para o Flow~C (somente modelo de linguagem) e a Tabela~\ref{tab:bugs-hybrid} para o Flow~B (híbrido). A coluna ``Aluc.'' indica o número de alucinações detectadas pela AST Validation, ou seja, hipóteses do modelo que referenciam construções inexistentes na AST real do programa.

\begin{table*}[htbp]
\caption{Detecção de bugs no Flow C (somente modelo de linguagem)}
\label{tab:bugs-llm}
\begin{center}
\begin{tabular}{lrrrrrrrrr}
\toprule
\textbf{Modelo} & \textbf{VP} & \textbf{FP} & \textbf{FN} & \textbf{P} & \textbf{R} & \textbf{F1} & \textbf{Acc} & \textbf{MCC} & \textbf{Aluc.} \\
\midrule
gpt-5.5               & 48 &  0 &  0 & 1{,}000 & 1{,}000 & 1{,}000 & 1{,}000 &  1{,}000 & 0  \\
claude-sonnet-4-6     & 48 &  4 &  0 & 0{,}923 & 1{,}000 & 0{,}960 & 0{,}982 &  0{,}938 & 1  \\
gemini-3.1-flash-lite & 46 &  3 &  2 & 0{,}939 & 0{,}958 & 0{,}949 & 0{,}946 &  0{,}825 & 1  \\
deepseek-r1:7b        & 27 &  8 & 21 & 0{,}771 & 0{,}563 & 0{,}651 & 0{,}654 &  0{,}386 & 7  \\
qwen2.5-coder:7b      & 34 & 20 & 14 & 0{,}630 & 0{,}708 & 0{,}667 & 0{,}636 & -0{,}059 & 13 \\
\bottomrule
\end{tabular}
\end{center}
\end{table*}

\begin{table*}[htbp]
\caption{Detecção de bugs no Flow B (Modelo de Linguagem + ESBMC)}
\label{tab:bugs-hybrid}
\begin{center}
\begin{tabular}{lrrrrrrrrrr}
\toprule
\textbf{Modelo} & \textbf{VP} & \textbf{FP} & \textbf{FN} & \textbf{P} & \textbf{R} & \textbf{F1} & \textbf{Acc} & \textbf{MCC} & \textbf{TCF} & \textbf{TRN} \\
\midrule
gpt-5.5               & 48 & 0 &  0 & 1{,}000 & 1{,}000 & 1{,}000 & 1{,}000 & 1{,}000 & 1{,}000 & N/A   \\
claude-sonnet-4-6     & 48 & 1 &  0 & 0{,}980 & 1{,}000 & 0{,}990 & 1{,}000 & 1{,}000 & 0{,}960 & 0{,}750 \\
gemini-3.1-flash-lite & 46 & 1 &  2 & 0{,}979 & 0{,}958 & 0{,}968 & 0{,}964 & 0{,}892 & 0{,}958 & 0{,}667 \\
deepseek-r1:7b        & 27 & 0 & 21 & 1{,}000 & 0{,}563 & 0{,}720 & 0{,}673 & 0{,}463 & 0{,}964 & 1{,}000 \\
qwen2.5-coder:7b      & 34 & 0 & 14 & 1{,}000 & 0{,}708 & 0{,}829 & 0{,}782 & 0{,}577 & 0{,}850 & 1{,}000 \\
\bottomrule
\end{tabular}
\end{center}
\end{table*}

\noindent\textit{Acc} e \textit{MCC}: calculados por função (positiva = tem bug, negativa = limpa). \textit{TCF}: proporção de hipóteses confirmadas pelo ESBMC. \textit{TRN}: proporção de falsos positivos do Flow~C eliminados no Flow~B. N/A: TRN indefinida pois o modelo já tinha FP\,=\,0 no Flow~C.

Os modelos de API superaram amplamente os modelos locais. O \texttt{gpt-5.5} foi o único a atingir desempenho perfeito (F1\,=\,1{,}000) já sem o ESBMC, sem gerar alarmes falsos nem perder nenhum bug real. O \texttt{claude-sonnet-4-6} encontrou todas as 48 instâncias de bug mas gerou 4 alarmes falsos no Flow~C; com o ESBMC, esse número caiu para 1, elevando o F1 para 0{,}990. O \texttt{gemini-3.1-flash-lite} perdeu 2 instâncias e gerou 3 FPs no Flow~C; o ESBMC reduziu para 1 FP (F1\,=\,0{,}968).

Os modelos locais tiveram desempenho inferior. O \texttt{deepseek-r1:7b} encontrou apenas 27 das 48 instâncias de bug (perdeu 21) e gerou 8 FPs. O ESBMC eliminou todos os 8 FPs, mas os 21 bugs perdidos são irrecuperáveis, pois o modelo simplesmente não os identificou. O \texttt{qwen2.5-coder:7b} encontrou 34 bugs mas gerou 20 FPs, o maior número de alarmes falsos; seu MCC negativo no Flow~C (${-}0{,}059$) indica desempenho abaixo do acaso nessa configuração. O ESBMC eliminou todos os 20 FPs, elevando a precisão de 0{,}630 para 1{,}000, mas os 14 bugs perdidos permaneceram sem detecção.

\subsection{Resultados por Categoria de Bug}

As Tabelas~\ref{tab:cat-llm} e~\ref{tab:cat-hybrid} detalham o desempenho por categoria de bug nos modos Flow~C e Flow~B, respectivamente. As categorias \texttt{assertion\_violation} e \texttt{division\_by\_zero} possuem 15 instâncias no \textit{ground truth}, enquanto \texttt{out\_of\_bounds} possui 18 instâncias distribuídas em 15 funções.

\begin{table*}[htbp]
\caption{Detecção por categoria de bug no Flow C (somente modelo de linguagem)}
\label{tab:cat-llm}
\begin{center}
\begin{tabular}{llrrrrr}
\toprule
\textbf{Categoria} & \textbf{Modelo} & \textbf{VP} & \textbf{FP} & \textbf{FN} & \textbf{P} & \textbf{F1} \\
\midrule
\multirow{5}{*}{\texttt{assertion\_violation}}
  & gpt-5.5               & 15 & 0 &  0 & 1{,}000 & 1{,}000 \\
  & claude-sonnet-4-6     & 15 & 0 &  0 & 1{,}000 & 1{,}000 \\
  & gemini-3.1-flash-lite & 13 & 0 &  2 & 1{,}000 & 0{,}929 \\
  & deepseek-r1:7b        & 11 & 1 &  4 & 0{,}917 & 0{,}815 \\
  & qwen2.5-coder:7b      & 14 & 3 &  1 & 0{,}824 & 0{,}875 \\
\midrule
\multirow{5}{*}{\texttt{division\_by\_zero}}
  & gpt-5.5               & 15 &  0 &  0 & 1{,}000 & 1{,}000 \\
  & claude-sonnet-4-6     & 15 &  4 &  0 & 0{,}789 & 0{,}882 \\
  & gemini-3.1-flash-lite & 15 &  3 &  0 & 0{,}833 & 0{,}909 \\
  & deepseek-r1:7b        &  9 &  4 &  6 & 0{,}692 & 0{,}643 \\
  & qwen2.5-coder:7b      & 15 &  8 &  0 & 0{,}652 & 0{,}789 \\
\midrule
\multirow{5}{*}{\texttt{out\_of\_bounds}}
  & gpt-5.5               & 18 &  0 &  0 & 1{,}000 & 1{,}000 \\
  & claude-sonnet-4-6     & 18 &  0 &  0 & 1{,}000 & 1{,}000 \\
  & gemini-3.1-flash-lite & 18 &  0 &  0 & 1{,}000 & 1{,}000 \\
  & deepseek-r1:7b        &  7 &  3 & 11 & 0{,}700 & 0{,}500 \\
  & qwen2.5-coder:7b      &  5 &  9 & 13 & 0{,}357 & 0{,}313 \\
\bottomrule
\end{tabular}
\end{center}
\end{table*}

\begin{table*}[htbp]
\caption{Detecção por categoria de bug no Flow B (Modelo de Linguagem + ESBMC)}
\label{tab:cat-hybrid}
\begin{center}
\begin{tabular}{llrrrrr}
\toprule
\textbf{Categoria} & \textbf{Modelo} & \textbf{VP} & \textbf{FP} & \textbf{FN} & \textbf{P} & \textbf{F1} \\
\midrule
\multirow{5}{*}{\texttt{assertion\_violation}}
  & gpt-5.5               & 15 & 0 &  0 & 1{,}000 & 1{,}000 \\
  & claude-sonnet-4-6     & 15 & 0 &  0 & 1{,}000 & 1{,}000 \\
  & gemini-3.1-flash-lite & 13 & 0 &  2 & 1{,}000 & 0{,}929 \\
  & deepseek-r1:7b        & 11 & 0 &  4 & 1{,}000 & 0{,}846 \\
  & qwen2.5-coder:7b      & 14 & 0 &  1 & 1{,}000 & 0{,}966 \\
\midrule
\multirow{5}{*}{\texttt{division\_by\_zero}}
  & gpt-5.5               & 15 & 0 &  0 & 1{,}000 & 1{,}000 \\
  & claude-sonnet-4-6     & 15 & 1 &  0 & 0{,}938 & 0{,}968 \\
  & gemini-3.1-flash-lite & 15 & 1 &  0 & 0{,}938 & 0{,}968 \\
  & deepseek-r1:7b        &  9 & 0 &  6 & 1{,}000 & 0{,}750 \\
  & qwen2.5-coder:7b      & 15 & 0 &  0 & 1{,}000 & 1{,}000 \\
\midrule
\multirow{5}{*}{\texttt{out\_of\_bounds}}
  & gpt-5.5               & 18 & 0 &  0 & 1{,}000 & 1{,}000 \\
  & claude-sonnet-4-6     & 18 & 0 &  0 & 1{,}000 & 1{,}000 \\
  & gemini-3.1-flash-lite & 18 & 0 &  0 & 1{,}000 & 1{,}000 \\
  & deepseek-r1:7b        &  7 & 0 & 11 & 1{,}000 & 0{,}560 \\
  & qwen2.5-coder:7b      &  5 & 0 & 13 & 1{,}000 & 0{,}435 \\
\bottomrule
\end{tabular}
\end{center}
\end{table*}

\textbf{assertion\_violation} é a categoria mais fácil. Todos os modelos de API acertam 100\% das acusações (P\,=\,1{,}000) já no Flow~C. Os modelos locais têm alguns erros de precisão no Flow~C, mas o ESBMC corrige: no Flow~B todos atingem P\,=\,1{,}000. Violações de asserção são detectáveis porque o código contém uma instrução explícita (\texttt{assert}) que o modelo consegue identificar visualmente.

\textbf{division\_by\_zero} tem desempenho mediano no Flow~C mas melhora bastante com o ESBMC. O \texttt{claude-sonnet-4-6} e o \texttt{gemini-3.1-flash-lite} encontram todos os 15 casos mas geram FPs que o ESBMC reduz. O \texttt{qwen2.5-coder:7b} também encontra todos os 15 mas gerou 8 FPs; o ESBMC eliminou todos, resultando em F1\,=\,1{,}000 no Flow~B.

\textbf{out\_of\_bounds} é a categoria mais difícil. Os três modelos de API encontram todos os 18 casos (R\,=\,1{,}000) tanto no Flow~C quanto no Flow~B. Já os modelos locais falham gravemente: o \texttt{deepseek-r1:7b} encontrou apenas 7 dos 18 e o \texttt{qwen2.5-coder:7b} apenas 5, bugs perdidos que o ESBMC não pode recuperar. Essa dificuldade é esperada: acessos fora dos limites dependem de raciocínio sobre índices e tamanhos, o que exige maior capacidade de compreensão do código.

\subsection{Resultados para Code Smells}

A Tabela~\ref{tab:smells} apresenta os resultados para detecção de \textit{code smells} (15 casos, sendo 5 por categoria). O Flow~B não aplica ESBMC a \textit{smells}, portanto os resultados são idênticos ao Flow~C para essa dimensão.

\begin{table*}[htbp]
\caption{Detecção de \textit{code smells} (Flow~C igual ao Flow~B)}
\label{tab:smells}
\begin{center}
\begin{tabular}{lrrrrrr}
\toprule
\textbf{Modelo} & \textbf{VP} & \textbf{FP} & \textbf{FN} & \textbf{P} & \textbf{R} & \textbf{F1} \\
\midrule
gpt-5.5               & 10 &  1 &  5 & 0{,}909 & 0{,}667 & 0{,}769 \\
claude-sonnet-4-6     & 10 &  0 &  5 & 1{,}000 & 0{,}667 & 0{,}800 \\
gemini-3.1-flash-lite &  5 &  0 & 10 & 1{,}000 & 0{,}333 & 0{,}500 \\
deepseek-r1:7b        &  0 &  0 & 15 & 0{,}000 & 0{,}000 & 0{,}000 \\
qwen2.5-coder:7b      & 15 & 95 &  0 & 0{,}136 & 1{,}000 & 0{,}240 \\
\bottomrule
\end{tabular}
\end{center}
\end{table*}

O \texttt{claude-sonnet-4-6} tem o melhor F1 para \textit{smells} (0{,}800), com precisão perfeita e R\,=\,0{,}667. O \texttt{gpt-5.5} fica logo atrás (F1\,=\,0{,}769), com 1 FP e mesma revocação. O \texttt{gemini-3.1-flash-lite} não emite FPs mas falha em detectar a maioria (R\,=\,0{,}333, F1\,=\,0{,}500). O \texttt{qwen2.5-coder:7b} detecta todos os 15 casos reais mas gera 95 FPs, tornando sua precisão inviável (P\,=\,0{,}136). O \texttt{deepseek-r1:7b} não detecta nenhum \textit{smell} (F1\,=\,0{,}000).

A Tabela~\ref{tab:smells-cat} detalha o desempenho por tipo de \textit{smell} para os três modelos de maior F1 em smells (\texttt{claude-sonnet-4-6}, \texttt{gpt-5.5} e \texttt{gemini-3.1-flash-lite}).

\begin{table*}[htbp]
\caption{Detecção por categoria de smell para os três modelos de maior desempenho}
\label{tab:smells-cat}
\begin{center}
\begin{tabular}{llrrrr}
\toprule
\textbf{Categoria} & \textbf{Modelo} & \textbf{VP} & \textbf{FP} & \textbf{FN} & \textbf{F1} \\
\midrule
\multirow{3}{*}{\texttt{complex\_conditional}}
  & claude-sonnet-4-6     & 5 & 0 & 0 & 1{,}000 \\
  & gpt-5.5               & 5 & 1 & 0 & 0{,}909 \\
  & gemini-3.1-flash-lite & 0 & 0 & 5 & 0{,}000 \\
\midrule
\multirow{3}{*}{\texttt{long\_method}}
  & claude-sonnet-4-6     & 0 & 0 & 5 & 0{,}000 \\
  & gpt-5.5               & 0 & 0 & 5 & 0{,}000 \\
  & gemini-3.1-flash-lite & 0 & 0 & 5 & 0{,}000 \\
\midrule
\multirow{3}{*}{\texttt{many\_parameters}}
  & claude-sonnet-4-6     & 5 & 0 & 0 & 1{,}000 \\
  & gpt-5.5               & 5 & 0 & 0 & 1{,}000 \\
  & gemini-3.1-flash-lite & 5 & 0 & 0 & 1{,}000 \\
\bottomrule
\end{tabular}
\end{center}
\end{table*}

Os três modelos detectam \texttt{many\_parameters} com precisão perfeita (F1\,=\,1{,}000). O \texttt{long\_method} é a categoria mais difícil: nenhum dos três detectou nenhum caso. Para \texttt{complex\_conditional}, o \texttt{claude-sonnet-4-6} obteve F1\,=\,1{,}000, o \texttt{gpt-5.5} ficou próximo (F1\,=\,0{,}909, com 1 FP), e o \texttt{gemini-3.1-flash-lite} não detectou nenhum caso.

\subsection{Discussão dos Resultados}

Os resultados confirmam quatro hipóteses centrais e revelam limitações que orientam trabalhos futuros.

\textbf{RQ1 — A verificação formal reduz os falsos positivos sem afetar a revocação.} Em todos os cinco modelos, o ESBMC reduziu FPs sem criar novos FNs. Isso é estrutural: o verificador só age sobre hipóteses que o modelo já reportou, então elimina FPs mas nunca recupera FNs. O ganho foi maior para os SLMs: \texttt{qwen2.5-coder:7b} eliminou 100\% de seus 20 FPs (TRN\,=\,1{,}000) e \texttt{deepseek-r1:7b} eliminou 100\% de seus 8 FPs. Para os modelos de API, o ganho foi menor: \texttt{claude-sonnet-4-6} reduziu FPs de 4 para 1 (TRN\,=\,0{,}750) e \texttt{gemini-3.1-flash-lite} de 3 para 1 (TRN\,=\,0{,}667). O \texttt{gpt-5.5} já não tinha FPs no Flow~C, então o TRN é indefinido. Isso corrobora Tihanyi et al.\ \cite{tihanyi2025}: o híbrido compensa especificamente a tendência a falsos positivos dos modelos.

\textbf{RQ2 — A taxa de alucinação separa modelos de API de modelos locais.} A AST Validation capturou muito mais alucinações nos SLMs locais: \texttt{deepseek-r1:7b} produziu 7 referências inválidas e \texttt{qwen2.5-coder:7b} produziu 13, contra 0 para \texttt{gpt-5.5} e 1 para cada um dos modelos de API restantes. Isso reflete diferenças de capacidade de raciocínio sobre código entre modelos de grande e pequeno porte, alinhadas com Ding et al.\ \cite{ding2025}. A AST Validation é especialmente crítica quando se usam SLMs locais.

\textbf{RQ3 — O ESBMC não compensa lacunas de revocação do modelo.} A diferença entre Flow~B e Flow~C é só de precisão: o ESBMC elimina FPs mas não recupera FNs. O \texttt{deepseek-r1:7b} ilustra esse limite: mesmo com TCF\,=\,0{,}964 e zero FPs no Flow~B, o F1 (0{,}720) fica abaixo dos modelos de API por causa dos 21 FNs irrecuperáveis. A categoria mais afetada é \texttt{out\_of\_bounds}: os SLMs locais detectam apenas 7 e 5 dos 18 casos (R\,=\,0{,}389 e 0{,}278), enquanto todos os modelos de API atingem R\,=\,1{,}000. Isso é consistente com Li et al.\ \cite{li2026}: modelos menores falham em bugs comportamentais complexos.

\textbf{RQ4 — \textit{Code smells} precisam de estratégias complementares.} O \texttt{many\_parameters} é detectado com precisão por todos os melhores modelos (F1\,=\,1{,}000), pois tem critério estrutural objetivo. Já o \texttt{long\_method} não foi detectado por nenhum dos três melhores modelos, sugerindo que funções longas mas coesas não ativam o padrão de \textit{smell} no \textit{prompting} usado. O \texttt{qwen2.5-coder:7b} tem \textit{over-reporting} severo (95 FPs em 15 casos reais, P\,=\,0{,}136), corroborando Souza et al.\ \cite{souza2026}. Como o ESBMC não cobre \textit{smells}, esses FPs não são corrigíveis, o que evidencia a necessidade de um filtro específico para \textit{smells} em versões futuras.

\textbf{Flow A como referência de validação do dataset.} O Flow~A confirmou todos os bugs com P\,=\,R\,=\,F1\,=\,1{,}000, o que era esperado: o dataset foi construído precisamente com o auxílio do ESBMC, e esse fluxo serve para atestar que os defeitos anotados são de fato disparáveis pelo verificador. A comparação relevante é entre Flow~B e Flow~C: o quanto a verificação formal agrega sobre a análise isolada do modelo.

\section{Metodologia}

\subsection{Etapas Metodológicas}

A pesquisa foi organizada nas etapas a seguir:

\begin{itemize}
    \item \textbf{Revisão da literatura} (concluída): cerca de 25 artigos em seis eixos: análise estática (SAST), verificação formal/BMC, modelos de linguagem para engenharia de software, detecção de bugs baseada em modelos de linguagem, abordagens híbridas e benchmarks.

    \item \textbf{Definição da arquitetura} (concluída): cinco etapas do pipeline, três modos de execução, categorias de bugs e smells, seis categorias de defeitos (três de bug e três de \textit{code smell}).

    \item \textbf{Implementação} (concluída): todas as etapas, incluindo a validação por tolerância da AST para filtrar alucinações em funções nativas do Python.

    \item \textbf{Curadoria do dataset} (concluída): construção manual de 70 funções Python distribuídas em seis categorias, sendo três de bugs executáveis (\texttt{division\_by\_zero}, \texttt{out\_of\_bounds} e \texttt{assertion\_violation}) e três de \textit{code smells} (\texttt{long\_method}, \texttt{many\_parameters} e \texttt{complex\_conditional}), além de 10 funções limpas como controle. O \textit{ground truth} foi anotado manualmente para cada função, indicando o tipo e a categoria do defeito presente. Cada função anotada como contendo um bug verificável foi submetida ao ESBMC-Python para confirmar que o defeito é disparável dentro do limite de execução configurado; todos os 45 casos foram confirmados com \texttt{VERIFICATION FAILED}.

    \item \textbf{Execução dos experimentos} (concluída): Flows A, B e C sobre o \textit{dataset} V1 com cinco modelos.

    \item \textbf{Análise dos resultados} (concluída): P/R/F1 por categoria e modelo, TCF e TRN.

    \item \textbf{Escrita e revisão do relatório} (concluída): redação do artigo com descrição do pipeline, análise dos resultados e comparação com trabalhos correlatos.
\end{itemize}



\section{Conclusão}

Este trabalho propôs e avaliou um pipeline híbrido que combina modelos de linguagem e o ESBMC-Python para detectar bugs executáveis e \textit{code smells} em programas Python.

Os resultados mostram que a confirmação formal reduz consistentemente os falsos positivos em todos os modelos. O \texttt{gpt-5.5} obteve F1 perfeito (1{,}000) já no modo somente-modelo. Entre os demais, o \texttt{claude-sonnet-4-6} teve o melhor F1 híbrido (0{,}990). O \texttt{qwen2.5-coder:7b} mostrou o maior ganho relativo de precisão com o ESBMC, passando de 0{,}630 para 1{,}000. Os resultados também confirmam que o ESBMC não recupera falsos negativos, e que \textit{code smells} precisam de filtros complementares.

Em relação ao EVA \cite{shivaji2026}, a proposta se diferencia em quatro pontos: verifica Python diretamente, sem traduzir para C; filtra alucinações via AST antes da verificação formal; trata \textit{code smells} separadamente; e usa seis categorias de classificação com métricas de P/R/F1 sobre um \textit{dataset} construído e validado manualmente.

Como próximos passos, pretende-se expandir o \textit{dataset} para uma versão V2, incluir avaliação estatística via \textit{bootstrap} e investigar estratégias de filtragem específicas para \textit{code smells}, com a possibilidade de usar os contraexemplos do ESBMC para validar o comportamento do código Python original.


\renewcommand{\refname}{Referências}
\bibliographystyle{IEEEtran}
\bibliography{references}

\end{document}